\documentclass[10pt,a4paper]{amsart}
\usepackage[margin=19mm]{geometry}
\usepackage{amsmath,amssymb,mathtools}
\usepackage[scr=rsfs,cal=euler]{mathalfa}
\usepackage{xfrac}
\usepackage[unicode,hidelinks,bookmarks=false]{hyperref}
\newcommand{\ie}{i.e.\ }
\newcommand{\cf}{cf.\ }
\newcommand{\resp}{respectively\ }
\newcommand{\Subsection}{\subsection}
\providecommand{\nobreakdash}{\nobreak\hbox{-}}
\setlength{\emergencystretch}{2em}
\multlinegap0pt
\usepackage{bm}
\usepackage{bbm}
\usepackage{dsfont}
\usepackage{xspace}
\usepackage{cancel}
\usepackage{mathtools}
\usepackage{amsthm}
\makeatletter
\@ifundefined{theorem}{\newtheorem{theorem}{Theorem}}{}
\makeatother
\title[Explicit Bicanonical Models of Eight Fake Quadrics]{Explicit Bicanonical Models of Eight Fake Quadrics}
\author{Lev Borisov, Carlos Rito}
\date{Source: arXiv:2608.12296v1 (CC BY 4.0)}
\begin{document}
\maketitle

\noindent\emph{Source and licence.} Explicit Bicanonical Models of Eight Fake Quadrics. Version arXiv:2608.12296v1 (CC BY 4.0), distributed under the Creative Commons Attribution 4.0 International licence, as recorded in the arXiv licence metadata for this exact version and shown on the abstract page https://arxiv.org/abs/2608.12296v1. Full rights notice: licenses/arxiv-2608.12296-CC-BY-4.0.txt.\par\smallskip

\noindent\textit{Editorial scope.} Counted unit: the Theorem whose source label is \texttt{None} in the article (source0.tex lines 704--706, source label \texttt{None}), delivered together with the article's own complete original proof of it (source0.tex lines 708--797, one \texttt{proof} environment of 90 lines). No mathematics is written, repaired or re-derived: statement and proof are the article's own text line for line. Required context retained uncounted: none. Every definition, display and label of those lines is retained unchanged. Omitted from this package and not redistributed: 1--703, 707--707, 798--830. Nothing inside a retained excerpt is omitted or altered: every retained body is delivered exactly as the article's lines are written, so no omission receipt of any kind accompanies this package. Citations: the retained excerpts carry no citation command at all, so no bibliography entry of the article is transcribed and no reference is redistributed. Reference adaptation: none was needed and none was made; every reference inside the delivered unit resolves inside it, so no unresolved reference or citation is produced. Printed slips kept literal: no printed word, symbol, hypothesis or number of the counted statement or of its proof was altered. Layout and class: the article's own class is \texttt{article} with packages including fontenc, lmodern, amsmath, amssymb, amsthm, amscd, mathtools, booktabs, geometry, microtype, hyperref; the wrapper is the lane's pinned \texttt{amsart} editorial wrapper at 10pt on A4, which restates the theorem-like environments the retained text needs and adds the article's own macro definitions that the retained text needs (none), with extra wrapper packages (bm, bbm, dsfont, xspace, cancel, mathtools). The delivered numbering is the unit's own. Assets: no retained span contains a figure environment, a graphics-inclusion command, a PDF-page inclusion, a table or a TikZ picture; no image file is copied into this package, the package declares no asset dependency and no image file is redistributed with it. Licence: version arXiv:2608.12296v1 of this article is distributed under the Creative Commons Attribution 4.0 International licence, as recorded in the arXiv licence metadata for this exact version and shown on the abstract page https://arxiv.org/abs/2608.12296v1. A full rights notice is delivered at licenses/arxiv-2608.12296-CC-BY-4.0.txt. Characters: every retained excerpt is pure ASCII, so the delivered engine, XeLaTeX with its default Latin Modern text font, reports no missing character.
\begin{theorem}
The eight fake quadrics given above are pairwise non-isomorphic.
\end{theorem}

\begin{proof}
For \(j=1,\ldots,8\), put
\[
        V_j=H^0(S_j,2K_{S_j}).
\]
By Riemann--Roch and Kodaira vanishing, \(\dim V_j=9\). The group
\(G\cong\mathbb Z/2\times\mathbb Z/4\) acts naturally on \(V_j\). For
every non-trivial element of order two, the periodicity in the
holomorphic Lefschetz formula gives trace $1$ on \(V_j\). An element of
order four fixes two points of \(S_j\). Each fixed point contributes
\[
        \frac{1}{(1-i)(1+i)}=\frac12,
\]
so its trace on \(V_j\) is again one. Hence the character of \(V_j\) is
the sum of the trivial and the regular characters of \(G\). In particular,
\(V_j\) has an eigenbasis with
\(\mathbb Z/2\times\mathbb Z/4\)-grading
\begin{equation}
 (0,0)\oplus(0,0)\oplus(0,1)\oplus(0,2)\oplus(0,3)
 \oplus(1,0)\oplus(1,1)\oplus(1,2)\oplus(1,3)
 \label{eq:bicanonical-grading}
\end{equation}
which we have also computed explicitly as the consequence of our construction of the covering surface.

Similarly, \(h^0(S_j,4K_{S_j})=49\), and the trace of every non-trivial
element of \(G\) on \(H^0(S_j,4K_{S_j})\) is one. Thus the invariant
eigenspace has dimension seven and each of the other seven eigenspaces
has dimension six. On the other hand,
\eqref{eq:bicanonical-grading} gives
\[
 \dim\!\left(\operatorname{Sym}^2V_j\right)_{(0,0)}=8,
 \qquad
 \dim\!\left(\operatorname{Sym}^2V_j\right)_{(0,2)}=7.
\]
Consequently, the kernel
\[
 I_2(S_j)=\ker\!\left(
   \operatorname{Sym}^2H^0(S_j,2K_{S_j})
   \longrightarrow H^0(S_j,4K_{S_j})
 \right)
\]
contains at least one quadric of character \((0,0)\) and one of
character \((0,2)\). For each \(j=1,\ldots,8\), the exact computation
gives
\(
        \dim I_2(S_j)=2.
\)

Choose a basis \(Q_{j,0},Q_{j,1}\) of \(I_2(S_j)\), and let
\(A_{j,0},A_{j,1}\) be the corresponding symmetric matrices. The
singular members of the pencil are parametrized by the vanishing of
\[
        \Delta_j(s,t)
        :=\det\bigl(sA_{j,0}+tA_{j,1}\bigr).
\]
If \(\Delta_j\) is non-zero, it is a binary form of degree \(9\), and
its zero divisor is the degeneracy divisor of the pencil. For one of
the eight surfaces, \(\Delta_j\) vanishes identically. In this case
every quadric in the pencil is singular and the degeneracy locus is the
whole of \(\mathbb P^1\).

An isomorphism between two of the surfaces would induce an isomorphism
between their bicanonical spaces, and hence a projective equivalence
between the corresponding degeneracy loci. Thus, for some
\[
        M=
        \begin{pmatrix}
        a&b\\ c&d
        \end{pmatrix}
        \in\operatorname{GL}_2(\overline{\mathbb Q}),
\]
the determinant forms corresponding to \(S_j\) and \(S_k\) would
satisfy
\[
        \Delta_j(as+bt,cs+dt)=\Delta_k(s,t),
\]
together with
\[
        1+n(ad-bc)=0
        \quad\text{for some }n,
\]
which imposes the condition \(\det(M)\neq0\).

For each of the \(28\) pairs \(1\leq j<k\leq8\), we formed the
corresponding scheme in the five-dimensional affine space with
coordinates \(a,b,c,d,n\). Exact computations in Magma show that all
these schemes are empty. Hence the eight degeneracy loci are pairwise
non-equivalent over \(\overline{\mathbb Q}\), and consequently the eight
fake quadrics are pairwise non-isomorphic.
\end{proof}

\end{document}
